%% =============================================================================
%% INTRODUCTION
%% =============================================================================

\section{Introduction} \label{sec:intro}

The statistical software landscape for econometrics has long been dominated by
\proglang{R} \citep{R}, \proglang{Stata}, and increasingly \proglang{Python}.
\proglang{R} in particular has developed a rich ecosystem for econometric analysis,
as documented by \cite{Zeileis+Koenker:2008} and more recently by
\cite{Michelaki+Tsagris:2025}, who provide a systematic overview of 207
econometrics-related \proglang{R} packages. The \pkg{AER} package
\citep{Kleiber+Zeileis:2008, AER} has become a standard reference for applied
econometrics in \proglang{R}. Alternative ecosystems have emerged in
\proglang{Julia} with \pkg{Econometrics.jl} \citep{Santiago-Calderon:2020}
and in \proglang{MATLAB} with dedicated toolboxes for panel data analysis
\citep{Alvarez+Barbero+Zofio:2017}. Comparative studies suggest that compiled
languages offer substantial performance advantages for computationally intensive
econometric applications \citep{Aruoba+Fernandez-Villaverde:2015,
Duarte+Duarte+Fonseca+Montecinos:2020}.

While these environments offer mature, well-tested implementations of econometric
methods, they present challenges for deployment in production systems, particularly
those requiring high performance, memory safety, or integration with modern AI
systems. Moreover, numerical accuracy remains a concern across statistical software
platforms, as documented in a series of papers by McCullough and colleagues
\citep{McCullough:1997, McCullough+Vinod:1999, McCullough:1999b}, who demonstrated
that users cannot safely assume their software is numerically accurate without
verification.

The \proglang{Rust} programming language has emerged as a compelling platform for
scientific computing, combining the performance characteristics of compiled
languages like \proglang{C++} with memory safety guarantees enforced at compile
time \citep{Rust}. The \proglang{Rust} data science ecosystem has matured rapidly,
with \pkg{polars} providing high-performance DataFrame operations \citep{polars},
Apache Arrow \pkg{DataFusion} enabling fast analytical query processing
\citep{Lamb+etal:2024}, and \pkg{linfa} \citep{linfa} offering machine learning
algorithms inspired by \proglang{Python}'s \pkg{scikit-learn}. However, a gap
remains: no comprehensive econometrics library exists for \proglang{Rust}.

\pkg{prompt2analytics} addresses this gap by providing a comprehensive econometrics
toolkit implemented in pure \proglang{Rust}. The software is built on \pkg{polars}
\citep{polars} for DataFrame operations, \pkg{ndarray} \citep{ndarray} for
n-dimensional arrays, \pkg{faer} \citep{faer} for linear algebra, and \pkg{statrs}
\citep{statrs} for statistical distributions.

The library provides a pure \proglang{Rust} implementation of regression, panel
data, instrumental variables, discrete choice, and time series methods.
Heteroskedasticity-consistent standard errors (HC0--HC3) follow the formulations
of \cite{White:1980} and \cite{MacKinnon+White:1985}, while high-dimensional
fixed effects estimation implements the alternating projections algorithm of
\cite{Gaure:2013}. A distinguishing feature is the integration with large
language models through the Model Context Protocol \citep{MCP:2024}, enabling
natural language interfaces to statistical analysis while keeping all data
processing local. The software is accessible through a command-line interface
and a desktop application. All implementations are validated against reference
\proglang{R} packages to ensure numerical accuracy.
